\section*{Chapitre \ref{ch:b3:stable-isotopes} --- Géochimie des isotopes stables}

\begin{solution}{exo:b3:stable-isotopes:1}
$\alpha = 1030.5/997.9 = 1.032669$; $1000\ln\alpha = 32.15$; $\Delta = 30.5 -
(-2.1) = \qty{32.6}{\permil}$. Elles diffèrent de \qty{0.45}{\permil}, soit le
$(\delta_A^2 - \delta_B^2)/2000 = 0.46$ de la proposition.
\end{solution}

\begin{solution}{exo:b3:stable-isotopes:2}
$\delta^{18}\mathrm O = 0.69\times(-25) - 13.6 = \qty{-30.9}{\permil}$;
$\delta\mathrm D = 5.6\times(-25) - 100 = \qty{-240}{\permil}$; $d = -240 +
8\times30.85 = \qty{6.8}{\permil}$.
\end{solution}

\begin{solution}{exo:b3:stable-isotopes:3}
$\Delta_B = -26.0 - 1.0 = \qty{-27.0}{\permil}$; $f_{\mathrm{org}} = (-5 -
1.0)/(-27.0) = 0.22$.
\end{solution}

\begin{solution}{exo:b3:stable-isotopes:4}
$T^2 = 0.0391\times10^6/(0.600 - 0.154) = \num{87668}$, $T = \qty{296.1}{K}$,
\qty{22.9}{\celsius}.
\end{solution}

\begin{solution}{exo:b3:stable-isotopes:5}
$f^{0.01} = 967/987$, $f = \exp(100\ln0.97974) = 0.129$~: \qty{87}{\percent} de
la vapeur a précipité. La pluie~: $1.0100\times967 - 1000 =
\qty{-23.3}{\permil}$.
\end{solution}

\begin{solution}{exo:b3:stable-isotopes:6}
$\alpha = 1000/960 = 1.041667$; $1000\ln\alpha = 40.82$; $1000(\alpha - 1) =
41.67$; $\Delta = 40$. Les deux approximations s'écartent d'environ \qty{0.8}{\permil}, la
seconde de deux fois davantage que la première~: à \qty{40}{\permil}, le développement
au premier ordre ne suffit plus.
\end{solution}

\begin{solution}{exo:b3:stable-isotopes:7}
$\delta_c = (1.50 + 29.99)/0.97001 = \qty{32.46}{\permil}$ VSMOW~; $1000\ln\alpha
= 1000\ln1.032464 = 31.95$; $T = 18030/(31.95 + 32.42) = \qty{280.1}{K}$,
\qty{7.0}{\celsius}. Le fractionnement varie d'environ \qty{0.23}{\permil} par
kelvin, donc une erreur de \qty{0.1}{\permil} sur l'eau vaut environ \qty{0.4}{\celsius}.
\end{solution}

\begin{solution}{exo:b3:stable-isotopes:8}
Effet de volume des glaces $125\times0.011 = \qty{1.375}{\permil}$~; le reste, $1.77 - 1.375 =
\qty{0.40}{\permil}$. Près de \qty{2}{\celsius}, la pente vaut $18030/275.15^2 =
\qty{0.238}{\permil.K^{-1}}$, donc l'eau profonde était environ \qty{1.7}{\celsius}
plus froide.
\end{solution}

\begin{solution}{exo:b3:stable-isotopes:9}
$\Delta^{33}\mathrm S = 8.2 - 0.515\times10.0 = \qty{3.05}{\permil}$, loin de 0~:
il s'agit de soufre à \omterm{def:b3:stable-isotopes:mif}{fractionnement indépendant de la masse}, qui ne se forme qu'en l'absence d'écran d'ozone. La pyrite
s'est formée sous une atmosphère dépourvue d'oxygène notable, avant environ
\qty{2.45}{\Ga}.
\end{solution}

\begin{solution}{exo:b3:stable-isotopes:10}
L'isotope lourd est conservé dans le système clos~: $R_0 = fR_v + (1 -
f)\alpha R_v$, donc $R_v/R_0 = 1/[f + \alpha(1 - f)]$. Pour $f = 0.2$, $\alpha =
1.01$~: $1/1.008$, la vapeur est plus légère de \qty{7.9}{\permil}~; la \omterm{prop:b3:stable-isotopes:rayleigh-distillation}{distillation de Rayleigh} donne
$0.2^{0.01}$, plus légère de \qty{16.0}{\permil}. Retirer chaque condensat aussitôt
pousse deux fois plus loin la vapeur.
\end{solution}

\begin{solution}{exo:b3:stable-isotopes:11}
$M = m(\delta_s - \Delta)/\Delta = 2000\times(-60 + 2.5)/(-2.5) = \qty{46000}{Gt}$.
Alors $m = 46000\times(-2.5)/(-25 + 2.5) = \qty{5100}{Gt}$, un peu en dessous des
valeurs publiées de 6000 à \qty{12000}{Gt}, dont la borne supérieure correspond à la
plus forte excursion.
\end{solution}

\begin{solution}{exo:b3:stable-isotopes:12}
$\delta_v = (1000 + \delta_{v,0})\exp\int(\alpha - 1)\,\mathrm d\ln f - 1000$, donc
$\mathrm d\delta_v = (1000 + \delta_v)(\alpha - 1)\,\mathrm d\ln f$, et $f =
e_s/e_{s,0}$ donne $\mathrm d\ln f = \mathrm d\ln e_s$. À \qty{-20}{\celsius},
$\alpha = 1.01872$~: $0.01872\times950.9\times0.0960 =
\qty{1.71}{\permil.K^{-1}}$ pour la vapeur.
\end{solution}

\begin{solution}{pb:b3:stable-isotopes:1}
\textbf{1.} $\alpha^{18}$~: 1.01605, 1.01691, 1.01779, 1.01872, 1.01968,
1.02068~; $\alpha^{\mathrm D}$~: 1.1424, 1.1524, 1.1630, 1.1744, 1.1866, 1.1997
(de $-5$ à \qty{-30}{\celsius}).

\textbf{2.} $980\times0.425^{0.01605} - 1000 = \qty{-33.4}{\permil}$.

\textbf{3.} $1.01691\times966.6 - 1000 = \qty{-17.0}{\permil}$.

\textbf{4.} Vapeur $966.6\times(0.3976)^{0.01779} - 1000 = \qty{-49.1}{\permil}$;
neige $1.01872\times950.9 - 1000 = \qty{-31.3}{\permil}$.

\textbf{5.} Vapeur $950.9\times(0.3680)^{0.01968} - 1000 = \qty{-67.6}{\permil}$;
neige $1.02068\times932.4 - 1000 = \qty{-48.3}{\permil}$.

\textbf{6.} $(932.4/980)\times0.0622 = 0.059$~: \qty{5.9}{\percent} de
l'oxygène-18, contre \qty{6.2}{\percent} de l'eau~: la vapeur a perdu l'isotope
lourd plus rapidement.

\textbf{7.} $(-17.0 + 48.3)/20 = \qty{1.57}{\permil}$ par degré.

\textbf{8.} $t = (-48.3 + 13.6)/0.69 = \qty{-50}{\celsius}$. La relation vaut pour
les températures de surface~; sur une calotte glaciaire, la surface est nettement plus froide que
l'air où la neige se forme, et le modèle néglige les effets de mélange et cinétiques.

\textbf{9.} Vapeur $-247.5$, $-352.5$, \qty{-462.7}{\permil}~; neige $\delta\mathrm D
= -132.8$, $-239.6$, \qty{-355.4}{\permil}.

\textbf{10.} $d = -132.8 + 8\times17.02 = 3.3$; puis 10.7 et
\qty{31.3}{\permil}~: l'excès croît vers l'intérieur des terres. La pente locale de la
distillation, $\mathrm d\delta\mathrm D/\mathrm d\delta^{18}\mathrm O = (\alpha^{\mathrm D} -
1)(1000 + \delta\mathrm D)/[(\alpha^{18} - 1)(1000 + \delta^{18}\mathrm O)]$, passe sous
8 à mesure que la vapeur devient très légère ($1000 + \delta\mathrm D$ diminue bien plus vite que
$1000 + \delta^{18}\mathrm O$), donc la neige s'écarte au-dessus de la droite.

\textbf{11.} La droite donne $8\times(-48.34) + 10 = \qty{-376.7}{\permil}$~; la
neige se place \qty{21}{\permil} au-dessus de celle-ci. Le modèle ne tient compte que
du fractionnement à l'équilibre~; la neige réelle se forme dans un air sursaturé,
avec un effet cinétique, et les masses d'air se mélangent à d'autres en cours de route.

\textbf{12.} $\delta_c = (3.25 + 29.99)/0.97001 = 34.27$; $1000\ln\alpha =
33.69$; $T = 18030/66.11 = \qty{272.7}{K}$, \qty{-0.4}{\celsius}.

\textbf{13.} $\delta_c = 36.09$, $1000\ln\alpha = 35.46$, $T = \qty{265.6}{K}$,
\qty{-7.5}{\celsius}~: plus froid que ne peut l'être l'eau de mer sans geler. Le
changement ne peut pas être entièrement dû à la température.

\textbf{14.} $\delta_w = 125\times0.011 = \qty{+1.375}{\permil}$; $1000\ln(1036.09/1001.375)
= 34.08$; $T = 18030/66.50 = \qty{271.1}{K}$, \qty{-2.0}{\celsius}~: un refroidissement de
\qty{1.6}{\celsius}.

\textbf{15.} $\Delta_{47} = 0.0391\times10^6/272.71^2 + 0.154 = 0.6797$ et
$0.0391\times10^6/271.12^2 + 0.154 = 0.6859$~: une différence de
\qty{0.006}{\permil}.

\textbf{16.} Elle mesure la façon dont les isotopes lourds sont répartis au sein du
minéral, un équilibre interne qui ne dépend que de la température.

\textbf{17.} $f_{\mathrm{org}} = (-5 - 0.5)/(-28) = 0.20$.

\textbf{18.} $m = 46000\times(-3.0)/(-60 + 3.0) = \qty{2400}{Gt}$.

\textbf{19.} $46000\times(-3.0)/(-25 + 3.0) = \qty{6300}{Gt}$.

\textbf{20.} $46000\times(-3.0)/(-5 + 3.0) = \qty{69000}{Gt}$, plus que le
réservoir lui-même, en un laps de temps géologiquement bref~: invraisemblable.

\textbf{21.} La profondeur et l'étendue de la dissolution des carbonates dans la carotte, qui
mesure la quantité de carbone entrée dans l'océan, indépendamment de sa composition~;
avec l'amplitude de l'excursion isotopique, elle fixe $\delta_s$.

\textbf{22.} La neige au Dome Ilse a $\delta^{18}\mathrm O =
\textbf{\qty[text-series-to-math=true]{-48.3}{\permil}} $ (\qty{-48.4}{\permil}
avec $\alpha$ variant de façon continue).
\end{solution}
